% !TeX root = ../main.tex

\chapter{Titolo del Capitolo}
\label{cap:nome_capitolo}

Qui inserisci l'introduzione al capitolo, spiegando brevemente di cosa parlerà questa parte della relazione.

% ---------------------------------------------------------
\section{Titolo della Sezione}
\label{sec:titolo_sezione}

Questo è un paragrafo di testo standard. Se devi citare un file di codice, puoi scriverlo usando la spaziatura fissa: \texttt{drone_simulator.py}. 

Se invece devi fare riferimento a un percorso di una cartella, puoi scriverlo così: \texttt{configs/drone.yaml}.

\subsection{Titolo della Sotto-sezione}
\label{subsec:titolo_sottosezione}

Ecco un esempio di elenco puntato per descrivere componenti o requisiti:
\begin{itemize}
    \item \textbf{Componente A:} Spiegazione del primo blocco.
    \item \textbf{Componente B:} Spiegazione del secondo blocco.
\end{itemize}

% ---------------------------------------------------------
\section{Analisi del Codice e Configurazioni}
\label{sec:codice}

Quando dovrete inserire i blocchi di codice interi (es. file YAML di Kubernetes o script Python), useremo un blocco come questo (che si appoggerà al pacchetto scelto nel main):

\begin{lstlisting}[language=Python, caption={Breve descrizione di cosa fa il codice}, label={lst:codice_esempio}]
# Questo è un commento nel codice Python
def setup_agent():
    print("Agent inizializzato")
\end{lstlisting}

Un blocco di codice simile può essere richiamato nel testo dicendo ad esempio: "come mostrato nel Listing~\ref{lst:codice_esempio}...".